\section{Dependence-aware HAC inference}\label{sec:method}

DA-HAC estimates the two components of~\eqref{eq:var} from the loss differentials of a single blocked cross-validation run. Figure~\ref{fig:architecture} gives an overview of the complete procedure: blocked cross-validation produces the loss differentials (stage~1); two parallel branches estimate the test-sampling variance, through variogram-calibrated kernel weights, and the training-induced variance, from the dispersion of fold means (stage~2); the two are combined into a variance, a reference distribution and diagnostics (stage~3). The steps are summarised in Algorithm~\ref{alg:dahac} and detailed below.

\begin{figure}[t]
\centering
\resizebox{\textwidth}{!}{% Architecture of the proposed comparison procedure (included by main.tex).
% Reading order: stage 1 left to right, stage 2 right to left, stage 3 left to right.
\begin{tikzpicture}[
    font=\footnotesize,
    >=Stealth,
    node distance=5mm and 5mm,
    box/.style={draw=#1!75!black, fill=#1!10, rounded corners=2pt, align=center,
                minimum height=13mm, text width=29mm, inner sep=2pt},
    box/.default=gray,
    lane/.style={draw=#1!70!black, dashed, rounded corners=4pt, inner sep=2.5mm},
    arr/.style={->, thick, draw=black!65},
    lbl/.style={font=\scriptsize\bfseries, text=#1!70!black, anchor=south west, inner sep=1pt},
]
\definecolor{cpipe}{HTML}{52514E}
\definecolor{ctest}{HTML}{2A78D6}
\definecolor{ctrain}{HTML}{EB6834}
\definecolor{cout}{HTML}{1BAF7A}

% ---------------- Stage 1: blocked cross-validation ----------------
\node[box=cpipe] (data) {\textbf{Data}\\[1pt] $(x_i,y_i)$ observed at\\ location $s_i$, time $t_i$};
\node[box=cpipe, right=of data] (block) {\textbf{Blocking}\\[1pt] $G$ space--time blocks\\ $\to K$ folds, buffer};
\node[box=cpipe, right=of block] (fit) {\textbf{Refit $A$ and $B$}\\[1pt] on each training set\\ $T_k$, predict fold $F_k$};
\node[box=cpipe, right=of fit] (diff) {\textbf{Loss differentials}\\[1pt] $d_i=\ell_A(i)-\ell_B(i)$\\ estimate $\hat\Delta=\bar d$};
\draw[arr] (data) -- (block);
\draw[arr] (block) -- (fit);
\draw[arr] (fit) -- (diff);
\begin{scope}[on background layer]
  \node[lane=cpipe, fit=(data)(diff)] (L1) {};
\end{scope}
\node[lbl=cpipe] at (L1.north west) {1\quad Blocked cross-validation (defines the estimand $\Delta_n$)};

% ---------------- Stage 2: DA-P variance (right to left) ----------------
\node[box=ctest, below=11mm of diff] (center) {\textbf{Fold-centring}\\[1pt] fold means $m_k$\\ $\tilde u_i=d_i-m_{k(i)}$};
\node[box=ctest, left=of center] (vario) {\textbf{Variograms of $\tilde u$}\\[1pt] in space and in time\\ $\to$ ranges $\hat r_S,\hat r_T$};
\node[box=ctest, left=of vario] (bw) {\textbf{Correlation model}\\[1pt] exponential + nugget\\ $\hat g_{ij}=\hat g_S(h_{ij})\,\hat g_T(\ell_{ij})$};
\node[box=ctest, left=of bw] (vu) {\textbf{Test-sampling var.}\\[1pt] $\hat V_u=\frac{\hat\sigma^2}{n^2}\sum_{i,j}\hat g_{ij}$\\ and per fold $\hat V_k$};
\draw[arr] (diff) -- (center);
\draw[arr] (center) -- (vario);
\draw[arr] (vario) -- (bw);
\draw[arr] (bw) -- (vu);

\node[box=ctrain, below=of center] (fm) {\textbf{Between-fold}\\[1pt] variance of the\\ fold means: $s_m^2$};
\node[box=ctrain, left=of fm, text width=63mm] (tau) {\textbf{Training variance} (method of moments)\\[1pt] $\hat\tau^2=\max\bigl\{0,\ s_m^2-K^{-1}\textstyle\sum_k\hat V_k\bigr\}$};
\node[box=ctrain, left=of tau] (vt) {\textbf{Nadeau--Bengio}\\[1pt] $\zeta_K\hat\tau^2$\\ $\zeta_K=\tfrac1K+\tfrac1{K-1}$};
\draw[arr] (center) -- (fm);
\draw[arr] (fm) -- (tau);
\draw[arr] (tau) -- (vt);
\draw[arr, dashed] ([xshift=-3mm]vu.south east) -- node[midway, below left=-1pt, font=\scriptsize]{$\hat V_k$} ([xshift=6mm]tau.north west);

\begin{scope}[on background layer]
  \node[lane=ctest, fit=(center)(vu)(fm)(vt)] (L2) {};
\end{scope}
\node[lbl=ctest] at (L2.north west) {2\quad DA-P variance estimation};

% ---------------- Stage 3: inference and diagnostics ----------------
\node[box=cout, below=11mm of vt.south west, anchor=north west, text width=40mm] (v) {\textbf{Total variance}\\[1pt] $\hat V=\hat V_u+\zeta_K\hat\tau^2$};
\node[box=cout, right=of v, text width=40mm] (df) {\textbf{Satterthwaite d.f.\ $\hat\nu$}\\[1pt] from $\hat n_{\mathrm{eff}}$, $\textstyle\prod_q 3L_q/(2\hat r_q)$\\ and $K-1$; range ratio $\hat r/L$};
\node[box=cout, right=of df, text width=40mm] (out) {\textbf{Output}\\[1pt] $T=\hat\Delta/\sqrt{\hat V}\sim t_{\hat\nu}$: $p$-value, CI\\ decision: reject / not / indet.};
\draw[arr] (vt.south) -- (vt.south |- v.north);
\draw[arr] (vu.west) -- ++(-5mm,0) |- (v.west);
\draw[arr] (v) -- (df);
\draw[arr] (df) -- (out);

\begin{scope}[on background layer]
  \node[lane=cout, fit=(v)(out)] (L3) {};
\end{scope}
\node[lbl=cout, anchor=south east] at (L3.north east) {3\quad Inference and diagnostics};
\end{tikzpicture}
}
\caption{Architecture of the proposed comparison procedure. Stage~1 (grey) is a standard blocked cross-validation that yields the per-observation loss differentials $d_i$ and defines the estimand. Stage~2 (DA-HAC) estimates the test-sampling variance $\hat V_u$ (blue branch, read from right to left: fold-centring, variograms, bandwidths, product-Bartlett kernel) and the training-induced variance $\zeta_K\hat\tau^2$ (orange branch). Stage~3 (green) combines them into the total variance, the Satterthwaite degrees of freedom and the test, together with the diagnostics used in Section~\ref{sec:discussion}.}\label{fig:architecture}
\end{figure}

\subsection{Fold-centred differentials}

Let $m_k$ be the mean of the $d_i$ in fold $k$ and $\tilde u_i=d_i-m_{k(i)}$. Under~\eqref{eq:model}, $\tilde u_i=u_i-\bar u_{k(i)}$ is free of the fold effects, so that the dependence structure of $u$ can be studied without confusion with the training-induced component.

\subsection{Variogram-calibrated bandwidths}

The key tuning parameter of any HAC estimator is the bandwidth, i.e.\ the distance beyond which covariances are neglected. Rules of thumb developed for time series~\cite{newey1987,andrews1991} do not transfer to space, and the dependence range of the loss differential is in general different from that of the response: it is the range of the \emph{residual} structure that both models fail to capture. We therefore estimate it from $\tilde u$ directly. The directional empirical semivariograms
\begin{align*}
\hat\gamma_S(h)&=\frac{1}{2|N_S(h)|}\sum_{(i,j)\in N_S(h)}(\tilde u_i-\tilde u_j)^2,& N_S(h)&=\{(i,j):t_i=t_j,\ \|s_i-s_j\|\in I_h\},\\
\hat\gamma_T(\ell)&=\frac{1}{2|N_T(\ell)|}\sum_{(i,j)\in N_T(\ell)}(\tilde u_i-\tilde u_j)^2,& N_T(\ell)&=\{(i,j):s_i=s_j,\ |t_i-t_j|\in I_\ell\},
\end{align*}
are computed over distance and lag bins $I_h$, $I_\ell$ up to a maximum lag (by default half the spatial extent and a quarter of the temporal extent). For purely spatial data, $N_S(h)$ uses all pairs. To each we fit the exponential model $\gamma(h)=c_0+c_1\{1-\exp(-h/a)\}$ by the weighted least squares of Cressie~\cite{cressie1985}, with weights $|N(h)|/\gamma(h)^2$. The \emph{practical range} $r=3a$ is the distance at which the correlation of the structured component falls to $5\%$; it is set to zero when the structured part explains less than $5\%$ of the sill ($c_1<0.05(c_0+c_1)$). The spatial and temporal bandwidths are
\[
b_S=\kappa\,\hat r_S,\qquad b_T=\kappa\,\hat r_T ,
\]
with $\kappa=1$ chosen in a pilot study on separate random seeds (Section~\ref{sec:sim1}). We also report the \emph{range ratio} $\max_q b_q/L_q$, the largest bandwidth relative to the extent of the domain along the same axis, as a diagnostic: when it approaches one, the sample contains only a handful of effectively independent regions, and no procedure can estimate the variance of $\hat\Delta$ reliably from it.

\subsection{Kernel estimate of the test-sampling variance}

We use the product-Bartlett kernel
\begin{equation}\label{eq:kernel}
w_{ij}=\prod_{q=1}^{p}\Bigl(1-\frac{|s_{iq}-s_{jq}|}{b_S}\Bigr)_+\Bigl(1-\frac{|t_i-t_j|}{b_T}\Bigr)_+ ,
\end{equation}
where a zero bandwidth keeps only pairs with identical coordinates on that axis. Each factor is a one-dimensional Bartlett (triangular) kernel, whose Fourier transform is non-negative; the product is therefore a positive semi-definite kernel on $\mathbb{R}^{p+1}$, which guarantees non-negative variance estimates. (A Bartlett kernel applied to the Euclidean distance in two or more dimensions does not have this property.) The test-sampling variance and its within-fold analogues are
\begin{equation}\label{eq:vu}
\hat V_u=\frac{1}{n^2}\sum_{i,j}w_{ij}\tilde u_i\tilde u_j,\qquad
\hat V_k=\frac{1}{n_k^2}\sum_{i,j\in F_k}w_{ij}\tilde u_i\tilde u_j .
\end{equation}
Only pairs with $w_{ij}>0$ contribute; they are enumerated with a $k$-d tree in the Chebyshev metric after scaling each coordinate by its bandwidth, at a cost proportional to the number of such pairs.

\subsection{Training-induced variance}

Under~\eqref{eq:model}, the between-fold variance of fold means exceeds their within-fold test-sampling variance by the variance of the fold effects. We estimate it by the method of moments,
\begin{equation}\label{eq:tau}
\hat\tau^2=\max\Bigl\{0,\ s_m^2-\frac1K\sum_{k=1}^K\hat V_k\Bigr\}.
\end{equation}
Since $\E[s_m^2]$ subtracts the covariance between fold effects (Proposition~\ref{prop:pairedt}), $\hat\tau^2$ targets $\tau^2(1-\rho)$ rather than $\tau^2$. The correlation $\rho$ is not identifiable from one cross-validation run~\cite{bengio2004}. For an assumed value $\rho_0$, the training component of~\eqref{eq:var} is $V_\eta=\tau^2[1/K+\rho_0(1-1/K)]=\zeta_K(\rho_0)\,\tau^2(1-\rho_0)$, with
\begin{equation}\label{eq:zeta}
\zeta_K(\rho_0)=\frac{1/K+\rho_0(1-1/K)}{1-\rho_0},\qquad \zeta_K(1/K)=\frac1K+\frac1{K-1}.
\end{equation}
Following Nadeau and Bengio~\cite{nadeau2003}, our default is the test fraction $\rho_0=n_{\text{test}}/n=1/K$, for which $\zeta_K$ is exactly their inflation $1/J+n_{\text{test}}/n_{\text{train}}$ with $J=K$ splits. The final variance estimate is
\begin{equation}\label{eq:vtotal}
\hat V=\hat V_u+\zeta_K(\rho_0)\,\hat\tau^2 .
\end{equation}

\subsection{What the variance estimator targets}

The following proposition makes explicit the conditions under which~\eqref{eq:vtotal} is unbiased, and the direction of its bias when they fail. We consider the untruncated estimator $\hat\tau^2_{\mathrm{raw}}=s_m^2-K^{-1}\sum_k\hat V_k$ and $\hat V_{\mathrm{raw}}=\hat V_u+\zeta_K(\rho_0)\hat\tau^2_{\mathrm{raw}}$.

\begin{proposition}\label{prop:bias}
Under the working model~\eqref{eq:model} with equal fold sizes, let $\bar m^u_k$ be the fold means of $u$, $v_u=K^{-1}\sum_k\Var(\bar m^u_k)$, $c_u=[K(K-1)]^{-1}\sum_{k\neq l}\Cov(\bar m^u_k,\bar m^u_l)$, and let $\beta_u=\E[\hat V_u]-V_u$ and $\bar\beta=K^{-1}\sum_k\{\E[\hat V_k]-\Var(\bar m^u_k)\}$ be the biases of the kernel estimators. Then
\begin{equation}\label{eq:bias}
\E[\hat V_{\mathrm{raw}}]-\Var(\hat\Delta)
=\underbrace{\beta_u-\zeta_K(\rho_0)\,(c_u+\bar\beta)}_{\text{test-sampling part}}
+\underbrace{\tau^2\Bigl[(1-\rho)\,\zeta_K(\rho_0)-\tfrac1K-\rho\bigl(1-\tfrac1K\bigr)\Bigr]}_{\text{training part}} .
\end{equation}
In particular, $\hat V_{\mathrm{raw}}$ is unbiased if (i) $\rho_0=\rho$, (ii) fold means of $u$ are uncorrelated across folds ($c_u=0$) and (iii) the kernel estimators are unbiased ($\beta_u=\bar\beta=0$). The training part is negative, i.e.\ the variance is underestimated, if and only if $\rho>\rho_0$.
\end{proposition}

\begin{proof}
The fold means are $m_k=\Delta_n+\eta_k+\bar m^u_k$ with $\eta$ independent of $u$, so their average variance is $\tau^2+v_u$ and their average covariance $\rho\tau^2+c_u$. By the identity in the proof of Proposition~\ref{prop:pairedt}, $\E[s_m^2]=(1-\rho)\tau^2+v_u-c_u$, and by definition $\E[K^{-1}\sum_k\hat V_k]=v_u+\bar\beta$, hence $\E[\hat\tau^2_{\mathrm{raw}}]=(1-\rho)\tau^2-c_u-\bar\beta$. With equal fold sizes, \eqref{eq:var} reads $\Var(\hat\Delta)=V_u+\tau^2[1/K+\rho(1-1/K)]$ and $\E[\hat V_u]=V_u+\beta_u$. Substituting gives~\eqref{eq:bias}. Writing $g(r)=1/K+r(1-1/K)$, the training part is $\tau^2[(1-\rho)g(\rho_0)/(1-\rho_0)-g(\rho)]$, which vanishes at $\rho=\rho_0$ and is strictly decreasing in $\rho$, with derivative $-\tau^2[g(\rho_0)/(1-\rho_0)+1-1/K]<0$.
\end{proof}

Proposition~\ref{prop:bias} separates three sources of error, each of which is evaluated empirically in Section~\ref{sec:sim1}. \emph{Training correlation.} For $K=5$ and $\rho_0=1/K$, the training part equals $\tau^2(0.25-1.25\rho)$: assuming $\rho_0=0.2$ when $\rho=0.5$ leaves $0.375\tau^2$ of the $0.6\tau^2$ training variance unaccounted for, whereas $\rho<\rho_0$ errs on the conservative side. The sensitivity of the test to this choice is reported in Section~\ref{sec:sim1rho}. \emph{Cross-fold covariance.} $c_u$ is the covariance between test-sampling fluctuations of different folds; it is reduced by allocating whole blocks to folds and by buffers, and vanishes when blocks are larger than the dependence range. Because it enters with a negative sign, residual cross-fold dependence makes the test anti-conservative. \emph{Kernel bias.} Fold-centring removes the fold-level average of $u$ from $\tilde u$, and the Bartlett kernel down-weights covariances near the bandwidth; both make $\beta_u$ and $\bar\beta$ negative when the dependence range is not small relative to the blocks and to the domain. The term $-\zeta_K\bar\beta$ partially offsets $\beta_u$: the part of the fold-level variance missed by the within-fold estimators is recovered by $\hat\tau^2$. Finally, truncation at zero, $\hat\tau^2=\max\{0,\hat\tau^2_{\mathrm{raw}}\}$, can only increase the variance, so $\E[\hat V]\ge\E[\hat V_{\mathrm{raw}}]$; its effect on the reference distribution is discussed next.

\subsection{Reference distribution}

HAC tests are known to over-reject when the effective sample size is small~\cite{kiefer2005}, and $\hat\tau^2$ rests on only $K-1$ degrees of freedom. We approximate the distribution of $\hat V$ by a scaled $\chi^2$ with Welch--Satterthwaite~\cite{satterthwaite1946} degrees of freedom
\begin{equation}\label{eq:df}
\hat\nu=\frac{\hat V^2}{\hat V_u^2/\nu_u+(\zeta_K\hat\tau^2)^2/(K-1)},\qquad
\nu_u=\min\Bigl\{\hat n_{\mathrm{eff}}-1,\ \prod_{q}\frac{3L_q}{2b_q}\Bigr\},
\end{equation}
with $\hat n_{\mathrm{eff}}=n^{-1}\sum_i\tilde u_i^2/\hat V_u$ truncated to $[2,n]$, where the product runs over the spatial and temporal axes with a positive bandwidth, and $L_q$ is the extent of the domain along axis $q$. The second term of $\nu_u$ is the classical approximation to the degrees of freedom of a kernel spectral estimate~\cite{andrews1991}: for a lag window $k$ of bandwidth $b$ on a record of length $L$, $\Var(\hat V)/V^2\approx2(b/L)\int k^2$, with $\int k^2=2/3$ for the Bartlett window, and the factors multiply for a product kernel. It caps the degrees of freedom when the bandwidth is not small compared with the domain, the regime in which HAC tests are most anti-conservative. The test statistic $T=\hat\Delta/\sqrt{\hat V}$ is referred to a Student $t_{\hat\nu}$ distribution, and the $(1-\alpha)$ confidence interval is $\hat\Delta\pm t_{\hat\nu,1-\alpha/2}\sqrt{\hat V}$.

When $\hat\tau^2_{\mathrm{raw}}\le0$ the training term vanishes and $\hat\nu=\nu_u$, which may be large even though only $K$ folds inform the training component. Truncation occurs when the observed dispersion of fold means is no larger than its predicted test-sampling part, that is, when the data show no evidence of a training effect; the variance itself is then not underestimated on average (Proposition~\ref{prop:bias}), but its sampling variability is understated. Section~\ref{sec:sim1} reports the rejection rate conditionally on truncation to check that this boundary case does not cause over-rejection.

\subsection{A three-way decision rule}\label{sec:rule}

Proposition~\ref{prop:bias} shows that no correction of the variance can be exact when the dependence range is comparable to the blocks and to the study domain: $c_u$, $\beta_u$ and $\bar\beta$ are then all non-negligible and are not identifiable from one sample. We therefore attach to DA-HAC an explicit decision rule with three outcomes:
\begin{equation}\label{eq:rule}
\text{decision}=\begin{cases}
\text{indeterminate} & \text{if } \max_q b_q/L_q>r^\ast,\\
\text{reject } H_0 & \text{if } \max_q b_q/L_q\le r^\ast\text{ and } p<\alpha,\\
\text{do not reject } H_0 & \text{otherwise,}
\end{cases}
\end{equation}
where $b_q/L_q$ is the estimated bandwidth relative to the extent of the domain along axis $q$ (the range ratio). An ``indeterminate'' outcome states that the sample contains too few effectively independent regions for the comparison to be resolved, whatever the observed difference. The threshold $r^\ast=\rstar$ was chosen in a pilot study on separate random seeds as the largest value for which the conditional rejection rate of a true null stayed within the Monte Carlo band of the nominal level, and was then validated on new seeds and on the other experiments (Section~\ref{sec:sim1rule}).

\begin{algorithm}[t]
\caption{DA-HAC test for $H_0:\Delta_n=0$}\label{alg:dahac}
\begin{algorithmic}[1]
\Require loss differentials $d_i$, folds $k(i)$, coordinates $(s_i,t_i)$, level $\alpha$, multiplier $\kappa$ (default 1)
\State $m_k\gets$ mean of $d_i$ over $F_k$; \ $\tilde u_i\gets d_i-m_{k(i)}$; \ $\hat\Delta\gets n^{-1}\sum_i d_i$
\State Fit exponential models to the spatial and temporal empirical variograms of $\tilde u$; \ $b_S\gets\kappa\hat r_S$, $b_T\gets\kappa\hat r_T$
\State Enumerate pairs with $w_{ij}>0$ under the product-Bartlett kernel~\eqref{eq:kernel}
\State Compute $\hat V_u$ and $\hat V_1,\dots,\hat V_K$ by~\eqref{eq:vu}
\State $\hat\tau^2\gets\max\{0,s_m^2-K^{-1}\sum_k\hat V_k\}$; \ $\hat V\gets\hat V_u+\zeta_K(\rho_0)\hat\tau^2$ \Comment{default $\rho_0=1/K$}
\State Compute $\hat n_{\mathrm{eff}}$, $\nu_u$ and $\hat\nu$ by~\eqref{eq:df}, and the range ratio $\max_q b_q/L_q$
\State $p\gets2\,\Pr(t_{\hat\nu}>|\hat\Delta|/\sqrt{\hat V})$; interval $\hat\Delta\pm t_{\hat\nu,1-\alpha/2}\sqrt{\hat V}$
\State \Return $p$, the interval, the diagnostics and the decision~\eqref{eq:rule}: indeterminate if $\max_q b_q/L_q>r^\ast$
\end{algorithmic}
\end{algorithm}

\subsection{Properties}

\begin{remark}[Special cases]
If $b_S=b_T=0$ and $\hat\tau^2=0$, then $\hat V=n^{-2}\sum_i\tilde u_i^2$ and DA-HAC reduces to the one-sample $t$-test on the $d_i$. If the data have no spatial (or no temporal) coordinate, the corresponding factor of~\eqref{eq:kernel} is omitted. Because the kernel is positive semi-definite, $\hat V\ge0$ always.
\end{remark}

\begin{remark}[Consistency]
For fixed bandwidths, $\hat V_u$ is a Conley-type spatial HAC estimator~\cite{conley1999}. Under increasing-domain asymptotics with mixing conditions and bandwidths that grow with, but more slowly than, the size of the domain, such estimators are consistent~\cite{conley1999,andrews1991}. Our bandwidths are data-dependent and our variance includes a component, $\hat\tau^2$, estimated from only $K$ folds; we therefore do not claim exact asymptotic validity, and instead evaluate the finite-sample behaviour of the whole procedure by simulation with known truth.
\end{remark}

\begin{remark}[Scope]
DA-HAC requires only the per-observation losses of the two models and their coordinates; it is agnostic to the learners, to the loss and to the blocking scheme. The same construction applies to any additive performance measure (squared error, absolute error, log loss, Brier score, 0--1 loss). Non-additive measures such as the AUC or the F1-score would require a linearisation (influence function) first.
\end{remark}
